\documentclass[11pt]{article}
\usepackage[a4paper,margin=25mm]{geometry}
\usepackage{amsmath,amssymb,amsthm,mathtools,microtype}
\usepackage[T1]{fontenc}\usepackage{lmodern}
\usepackage[colorlinks=true,linkcolor=blue,urlcolor=blue]{hyperref}
\newtheorem{theorem}{Theorem}[section]
\newtheorem{lemma}[theorem]{Lemma}
\newcommand{\F}{\mathbb F}
\newcommand{\SV}{\operatorname{SV}}
\title{Polarization and Degree-Independent Direction Grids\\
for Rank-One or Invertible Affine Words}
\author{icekylinx \and GPT-6 Astra Pro}\date{26 October 2026}
\begin{document}\maketitle
\begin{abstract}
For local ranks one or \(w\), the constant-rowspace concentration
theorem gives an affine-orbit hitting map with \(2w-1\) seeds in
characteristic zero or characteristic greater than the local
entry-degree bound. This note chooses its direction grids
independently of that degree.
The lowest homogeneous term of each scalar numerator has individual
degrees below the characteristic, so its full polarization is nonzero
even when its total degree is at least the characteristic.
The coefficient is governed by a product of individual factorials,
not the factorial of the total degree.
For fixed \(w\), this gives
\(O_w(n^{2w-1}D^w)=O_w(n^{3w-1}\delta^w)\) oracle queries,
where \(D=n\delta\).
\end{abstract}

\noindent\textbf{Status.}
This quantitative argument has independent model-assisted internal
reconstructions and an exact-source audit. This is not external
peer review, and no publication-priority claim is made.

\section{Setting and theorem}

Let \(n,w,\delta\ge1\).
Assume \(\operatorname{char}\F=0\) or
\(\operatorname{char}\F=p>\delta\).
Consider a nonzero polynomial
\[
 f(X)=uA_1(\ell_1(X)+c_1)\cdots A_n(\ell_n(X)+c_n)v,
\]
where \(u,v\) are constant endpoints, \(L=(\ell_1,\ldots,\ell_n)\)
is invertible and linear, and each \(w\)-by-\(w\) polynomial matrix
\(A_i(T)\) has entry degree at most \(\delta\) and generic rank
one or \(w\).
The degree bound is on a representation with this rank pattern;
a scalar individual-degree promise alone does not suffice.

We use \emph{Constant Row Spaces and Low-Degree Concentration for
Invertible Affine ROABP Words}, in
\path{invertible_affine_constant_space_concentration.tex}.
Its structural theorem is the only dependency.
Write \(D=n\delta\), fix \(n\) distinct interpolation abscissas,
and let \(\lambda_i(y)\) be the Lagrange basis. Put \(k=w-1\) and
\[
 \mathcal H_i(t,\boldsymbol y,\boldsymbol z)
 =t^i+\sum_{j=1}^{k}z_j\lambda_i(y_j).
\]

\begin{theorem}\label{thm:grid}
For each \(1\le j\le k\), choose arbitrary sets \(Y_j,Z_j\subseteq\F\)
with
\[
 |Y_j|=n^2,\qquad |Z_j|=D+1.
\]
Choose a set \(T\subseteq\F\) with \(|T|=nD+1\).
Elements within each set are distinct; the sets may overlap.
The Cartesian evaluations of \(\mathcal H\) hit every nonzero
polynomial in the stated class.
The number of oracle queries is at most
\begin{equation}\label{eq:count}
 n^{2w-2}(nD+1)(D+1)^{w-1}
 =O_w(n^{2w-1}D^w)=O_w(n^{3w-1}\delta^w).
\end{equation}
A working field of at least \(nD+1\) elements suffices.
\end{theorem}

For \(w=1\) there are no direction or amplitude parameters and
the displayed product of their grids is empty.
If \(n=0\) or \(\delta=0\), one evaluation handles the constant case.
Over a smaller field one may pass to a sufficiently large extension,
provided oracle evaluations there are available.

\section{Polarization without a total factorial}

\begin{lemma}\label{lem:polar}
Let \(K\) be any field, and let \(0\ne H(Y)\in K[Y_1,\ldots,Y_n]\)
be homogeneous of degree \(m\).
Assume characteristic zero, or that every exponent of every
nonzero monomial of \(H\) is strictly less than
\(p=\operatorname{char}K\).
For \(m\) independent vector variables \(V_1,\ldots,V_m\), define
\[
 \mathcal P_H(V_1,\ldots,V_m)
 =[z_1\cdots z_m]\,
 H\!\left(\sum_{j=1}^m z_j V_j\right).
\]
Then \(\mathcal P_H\ne0\) and is multilinear in the vector arguments.
For any invertible linear map \(L\),
\[
 \mathcal P_{H\circ L}(V_1,\ldots,V_m)
 =\mathcal P_H(LV_1,\ldots,LV_m)\ne0.
\]
\end{lemma}
\begin{proof}
Write \(H=\sum_{|\alpha|=m}h_\alpha Y^\alpha\).
Fix an ordered tuple \(i_1,\ldots,i_m\) in which index \(i\)
occurs \(\alpha_i\) times. The coefficient of
\[
 \prod_{j=1}^m (V_j)_{i_j}
\]
in \(\mathcal P_H\) is exactly
\[
 h_\alpha\prod_{i=1}^n\alpha_i!.
\]
Indeed each group of \(\alpha_i\) distinct amplitude variables
assigned to \(Y_i^{\alpha_i}\) contributes its own
multinomial factor \(\alpha_i!\).
The ordered vector monomial determines \(\alpha\), so different
input monomials do not cancel this coefficient.

For a nonzero \(h_\alpha\), every \(\alpha_i!\) is nonzero under
the hypothesis. Hence \(\mathcal P_H\ne0\).
This calculation neither divides by \(m!\) nor assumes \(m<p\).
The covariance identity follows directly by substitution, and
invertibility in every vector argument preserves nonzeroness.
\end{proof}

\begin{lemma}\label{lem:directions}
Under the hypotheses of Lemma~\ref{lem:polar}, the polynomial
\[
 \mathcal P_H(L\lambda(y_1),\ldots,L\lambda(y_m))
\]
is nonzero and has degree at most \(n-1\) in each direction
variable, where \(\lambda(y)=(\lambda_1(y),\ldots,\lambda_n(y))\).
The assertion holds over any coefficient-field extension.
\end{lemma}
\begin{proof}
The nonzero polarized form after \(L\) is a tensor in \(m\)
independent vector arguments. Substituting the independent basis
polynomials \(\lambda_i(y_j)\) in the \(j\)-th argument is an
injective linear map for each argument. Their tensor product
is injective, so the result stays nonzero.
Multilinearity gives the individual degree bound \(n-1\).
\end{proof}

The transformed polynomial \(H\circ L\) may have exponents at least
\(p\). This does not invalidate the argument: the nonzero polarization
is established before the substitution and then transported by the
displayed covariance identity.

\section{Apply the tensor coefficient to every scalar segment}

As in the structural theorem, split each rank-one local factor
by a nonzero pivot:
\[
 R(T)=(R(T)e_b)\bigl(e_a^{\mathsf T}R(T)/h(T)\bigr),
 \qquad h=R_{ab}\ne0.
\]
Keep the scalar denominator \(h\) outside the numerator word.
If there are \(r\le n\) rank-one factors, this gives \(r+1\le n+1\)
nonzero scalar numerators, together with nonzero univariate
denominators.

Each numerator has an invertible middle run and polynomial
univariate endpoint vectors of degree at most \(\delta\).
Completing these endpoints with constant standard basis vectors
preserves the degree bound.
Every original variable occurs once within a numerator, though
neighboring numerators may share boundary variables.
Consequently each numerator has individual latent degree at most
\(\delta\).

Set \(\gamma_i(t)=t^i\) and
\(C_i(t)=c_i+\ell_i(\gamma(t))\).
Every \(C_i\) is transcendental over \(\F\).
Work over \(E=\F(t)\).
The structural concentration theorem gives, for each numerator
vanishing at \(C\), a lowest nonzero homogeneous residual term
\[
 H_\nu(Y)\quad\text{of degree }1\le m_\nu\le k=w-1.
\]
Its individual degrees remain at most \(\delta\).
In characteristic \(p>\delta\), these are all strictly below \(p\).
Lemma~\ref{lem:polar} therefore gives a nonzero full polarization.

After the original-coordinate substitution
\[
 Y=L\left(\sum_{j=1}^{k}z_j\lambda(y_j)\right),
\]
the coefficient of \(z_1\cdots z_{m_\nu}\), with every other
amplitude exponent zero, is the nonzero polynomial
\[
 P_\nu(y_1,\ldots,y_k)
 =\mathcal P_{H_\nu}
       (L\lambda(y_1),\ldots,L\lambda(y_{m_\nu})).
\]
It does not depend on unused direction variables.
Every higher homogeneous residual term has larger total amplitude
degree, so it cannot alter this coefficient.
By Lemma~\ref{lem:directions},
\[
 \deg_{y_j}P_\nu\le n-1
 \quad\text{for every }j.
\]
If a numerator is already nonzero at \(C\), use its nonzero
constant amplitude coefficient for \(P_\nu\).

The product \(P=\prod_\nu P_\nu\) is nonzero over \(E\), and
\[
 \deg_{y_j}P\le(n+1)(n-1)=n^2-1
 \quad\text{for every direction }j.
\]
Thus every Cartesian product of \(n^2\)-point direction grids
contains a tuple for which \(P\ne0\).
At that tuple every scalar numerator stays nonzero as a rational
function of \(t,\boldsymbol z\).
All pivot denominators stay nonzero as well: their constant
amplitude terms are the nonzero \(h(C_i(t))\), independently
of the chosen directions.
Their exact product identity proves that the original polynomial
is nonzero after those direction choices.

\section{Final evaluation count}

\begin{proof}[Proof of Theorem~\ref{thm:grid}]
Choose the direction tuple established above.
The original polynomial \(f(\mathcal H)\) is then nonzero in
\(t,z_1,\ldots,z_k\).
Its individual degrees are at most \(nD\) in \(t\) and at most
\(D\) in every amplitude. Therefore it cannot vanish on
\(T\times Z_1\times\cdots\times Z_k\).
The full Cartesian search includes that tuple and hence hits \(f\).
Counting all direction and remaining seed values gives
\eqref{eq:count}.

Since \(D=n\delta\ge n\), we have \(nD+1\ge n^2+1\), as well
as \(D+1\) and \(n\). This field size supplies all grids and
interpolation nodes.
The actual queries evaluate the original polynomial; they need
not separately avoid numerical zeros of auxiliary pivot
denominators used only in the symbolic proof.
\end{proof}

This addendum changes only the finite grid. The characteristic
hypothesis and local-rank promise are unchanged.
It does not convert latent digit-weight concentration into an
all-characteristic affine hitting theorem.
\end{document}
